\paragraph{Data.} The benchmark uses functional T-cell readouts measured by
interferon-$\gamma$ ELISpot on neoantigen candidates. All peptide sequences,
HLA assignments, and ELISpot spot-forming-count values were de-identified before
analysis and carry no patient-identifying information. To preserve the
double-blind review process, the originating cohort and any contributing
institutions are withheld from this version and will be named, with the
appropriate accession or controlled-access details, in the camera-ready version.

\paragraph{Code.} The harmonization pipeline, the per-tool scoring wrappers, the
per-patient evaluation protocol, and the analysis scripts that produce every
table and figure in this paper are organized for reproducible re-execution and
will be released in a public repository upon publication. To support peer review,
an anonymized snapshot of the code together with the harmonized, de-identified
derived score tables (per tool, per peptide) and the per-patient correlation
outputs underlying every table and figure are made available to reviewers and
editors, so that all reported metrics can be regenerated end to end without access
to any patient-identifying data.

\paragraph{Restricted tools and pending results.} A subset of predictors is
distributed under academic data-use agreements that restrict redistribution and
the publication of derived numbers---these include the netMHCpan, netMHCstabpan,
and NetTepi families and the related expansion-wave tools integrated into the
pipeline. The results for these tools are \emph{not} included in the main
conclusions of this paper and are marked \pendingDTU\ throughout. They will be
reported only after the required written data-use consent has been obtained, and
the full provenance and licensing status of every tool is tracked separately
(see the accompanying provenance record). The headline claims of this paper do
not depend on any restricted tool.
